\documentclass[11pt]{article}
\usepackage{../../../tex/lecturestyle}

\title{Persistent Segment Tree}
\author{}
\date{}

\begin{document}
\thispagestyle{fancy}

\lecturetitleband{Data Structures}{Persistent Segment Tree}{7}

\section{Overview}
An ordinary data structure is \emph{ephemeral}: an update discards the old
state and keeps only the new one. A \emph{persistent} data structure keeps
every version reachable instead, and a \emph{fully persistent} one --- the
kind this lecture builds --- allows querying \emph{and further updating}
any past version, not just the most recent one. The \textbf{persistent
segment tree} is the most broadly useful persistent structure in
competitive programming: cheap to build from an ordinary pointer-based
segment tree, and the basis for a wide family of range-query techniques.
Assumes familiarity with segment trees and lazy (range-update) segment
trees.

\subsection{Why this is needed}
\begin{definitionbox}
Array $A_0 = (A_{0,1}, \ldots, A_{0,N})$ is given. Each query either sets
$A_t$'s $i$-th entry to $x$ (creating a new array $A_q$), or asks for the
minimum of $A_t$ over a range --- for an arbitrary earlier version $t$, not
necessarily the most recent one.
\end{definitionbox}

If every query is known in advance (\emph{offline}), this is manageable:
treat each update as spawning a new version from an existing one, so the
versions form a tree, and answer queries while DFS-ing that tree with an
ordinary segment tree, undoing each update on the way back up. But if
queries must be answered \emph{online} --- each one before the next is
revealed --- that rollback trick is unusable, since which version comes
next isn't known ahead of time. Storing a full copy of the array at every
version costs $O(N)$ time and space per update, which is usually too slow.
A persistent segment tree answers every such query in $O(\log N)$ time and
space, online, by sharing almost everything between versions.

\section{Structure}

\subsection{Pointer-based segment trees, briefly}
A persistent segment tree is built on top of a \emph{pointer-based} (or
\emph{recursive}) segment tree: an internal node holds two pointers, to
its left and right child, rather than living at a fixed array index. An
ordinary point update on such a tree walks root to leaf and overwrites the
$O(\log N)$ nodes on that path; the pointer structure elsewhere is
untouched.

\subsection{Path copying: element updates}
The problem with overwriting those $O(\log N)$ nodes is that it destroys
the old version. The fix is \textbf{path copying}: instead of mutating
those nodes in place, clone each one, write the updated value into the
clone, and point the clone's children at whichever \emph{child} nodes are
correct for the new version (freshly cloned children where the path
continues, the original untouched children everywhere else). The clones
form a new root-to-leaf path; every other pointer in the entire tree
--- and so every subtree not on that path --- is simply shared between
old and new versions, unmodified.

\begin{theorembox}
A persistent segment tree update touches $O(\log N)$ nodes (one clone per
level) and creates a new, independent root. Both old and new versions
remain fully valid and queryable, and the new version costs
$O(\log N)$ additional time and space over the old one --- not $O(N)$.
\end{theorembox}

Versions are identified purely by their root: distinct versions have
distinct roots, and querying is done by walking down from whichever root
names the desired version, exactly as in an ordinary segment tree.

% FIGURE: pst_pathcopy
\begin{figure}[H]
\centering
\begin{tikzpicture}[x=1.15cm,y=1.15cm,
  old/.style={lvert,minimum size=8.5mm,font=\scriptsize\sffamily,fill=lectureink!8},
  new/.style={lvert,minimum size=8.5mm,font=\scriptsize\sffamily,fill=lectureaccent!30,draw=lectureaccent,thick}]
\node[old] (R0) at (2,3) {$[0,4)$};
\node[old] (A)  at (1,1.5) {$[0,2)$};
\node[old] (B)  at (3,1.5) {$[2,4)$};
\node[old] (L0) at (0.5,0) {$0$};
\node[old] (L1) at (1.5,0) {$1$};
\node[old] (L2) at (2.5,0) {$2$};
\node[old] (L3) at (3.5,0) {$3$};
\node[new] (R1) at (6.2,3) {$[0,4)'$};
\node[new] (B1) at (5.6,1.5) {$[2,4)'$};
\node[new] (L3n) at (6.4,0) {$3'$};
\foreach \p/\c in {R0/A,R0/B,A/L0,A/L1,B/L2,B/L3}{ \draw[lthin] (\p) -- (\c); }
\draw[lthin,lectureaccent,thick] (R1) -- (B1);
\draw[lthin,lectureaccent,thick] (B1) -- (L3n);
\draw[lthin,lectureaccent,thick,dashed] (R1) to[bend right=14] (A);
\draw[lthin,lectureaccent,thick,dashed] (B1) -- (L2);
\node[llabel,anchor=west] at (7.2,3) {version 1 root};
\node[llabel,anchor=east] at (1.0,3) {version 0 root};
\end{tikzpicture}
\caption{Path copying. Version 0 is the full tree on positions $0,\dots,3$. Setting position 3 creates only the new nodes on the root-to-leaf path (shaded, primed): a new root, a new right child and a new leaf. Every other node, namely the left subtree $[0,2)$ and leaf $2$, is shared with version 0 through the dashed pointers, so the update costs $O(\log n)$ time and memory and version 0 remains intact.}\label{fig:pst_pathcopy}
\end{figure}

\subsection{Path copying: range updates}
The same idea persists a \emph{lazy} segment tree. A lazy range update
proceeds in three passes: push pending lazy values down to children along
the way to the target range, apply the update operator at the $O(\log N)$
nodes exactly covering the range, and recombine values back up to the
root. Only $O(\log N)$ nodes ever change value across all three passes, so
the same rule applies: whenever a node's value or lazy tag is about to
change, clone it first (and, when \emph{pushing} a lazy tag down, clone
the children being pushed into as well, then push into the clones ---
never into the shared originals).

\subsection{Copying an identical range between versions}
A capability an ordinary segment tree has no analogue for: given two
versions, copy one version's data over some range $[l, r]$ into the other,
producing a third version. Since $[l, r]$ decomposes into $O(\log N)$
subtree-aligned pieces (the same decomposition an ordinary range query
uses), this is just $O(\log N)$ pointer relinks --- graft the source
version's node for each of those pieces directly into the destination
version's copied path, with no need to touch what's inside those
subtrees at all.

\section{Worked Applications}

\subsection{A chain of frequency arrays}
A recurring and very productive pattern: build a \emph{chain} of versions
$S_0, S_1, \ldots, S_N$ (rather than a branching tree of them) over a
value domain, where $S_i$ is obtained from $S_{i-1}$ by incrementing the
count at value $A_i$. Then $S_i$'s prefix sum up to some value $v$ is
exactly the count of $v' \le v$ among $A_1, \ldots, A_i$ --- the frequency
array evolves one insertion at a time, and every intermediate state stays
queryable.

\begin{examplebox}[: range count $\le x$]
The count of values $\le x$ in $A_l, \ldots, A_r$ is (prefix sum of $S_r$
up to $x$) $-$ (prefix sum of $S_{l-1}$ up to $x$), a single query against
each of two versions --- $O(\log N)$ total, after $O(N \log N)$
preprocessing to build the chain.
\end{examplebox}

\begin{examplebox}[: range $k$-th smallest]
The $k$-th smallest value in $A_l, \ldots, A_r$ is the least $v$ such that
(prefix count of $S_r$ up to $v$) $-$ (prefix count of $S_{l-1}$ up to $v$)
is at least $k$. Binary searching $v$ with an $O(\log N)$ query at each
step gives $O(\log^2 N)$, but there is a faster way: since $S_r$ and
$S_{l-1}$ have \emph{identical shape} (same domain, built by the same
sequence of splits), descend both trees \emph{simultaneously}. At each
pair of corresponding nodes, the difference between their left children's
counts tells you how many of the $k$ smallest candidates lie in the left
half; branch into left or right accordingly (adjusting $k$ if going
right), on both trees at once. This is the segment-tree ``binary search on
the tree'' technique applied to two versions together, and it costs only
$O(\log N)$.
\end{examplebox}

\begin{notebox}
The same two techniques answer range count/order-statistic queries on tree
paths (not just array ranges): give each vertex $v$ its own persistent
frequency tree, extending its parent's version by one insertion (its own
value). A root-to-$v$ query, combined by inclusion–exclusion with a
root-to-$w$ query and the shared root-to-$\mathrm{lca}(v,w)$ query
(subtracted twice, since it is double-counted, with $\mathrm{lca}(v,w)$'s
own value added back in), reduces a path query on the tree to the exact
same simultaneous-descent trick above. Building all $N$ per-vertex trees
this way, each extending its parent by a single insertion, costs
$O(N \log N)$ total.
\end{notebox}

\subsection{Beyond array and tree-path queries}
The same-range-copy operation opens up applications with no obvious
connection to persistence at first glance: some counting or DP problems
turn out to satisfy a recurrence where one state is built out of two
\emph{other} states by copying disjoint sub-ranges from each into a fresh
array, which a persistent (lazy) segment tree implements directly in
$O(\log N)$ per copy rather than reconstructing the state from scratch.
Recognizing this pattern is the harder part; working through a full
example is left to the further-reading problems below.

\section{Complexity}
\begin{center}
\begin{tabular}{ll}
\toprule
\textbf{Operation} & \textbf{Time / extra space} \\
\midrule
Build initial version & $O(N)$ \\
Point update (new version) & $O(\log N)$ \\
Range update, lazy (new version) & $O(\log N)$ \\
Range query on any version & $O(\log N)$ \\
Copy a range between two versions & $O(\log N)$ \\
Range $k$-th smallest via chained versions & $O(\log N)$ per query, $O(N \log N)$ build \\
\bottomrule
\end{tabular}
\end{center}

\section{Reference Implementations}
Both examples below implement two variants: a branching multi-version
array with point update and range-min query, and a chained frequency-tree
supporting range count-$\le x$ and range $k$-th-smallest via the
simultaneous-descent trick.

\codenote
\subsection{Rust}
\corelisting{Rust}{example.rs}

\subsection{Python}
\corelisting{Python}{example.py}

\section{Worked Example}
\begin{examplebox}
Starting from version $v_0 = (5,9,4,7)$, updating position $2$ to $2$
gives $v_1 = (5,2,4,7)$, while updating position $4$ of $v_0$ (not $v_1$!)
to $1$ gives a separately branching $v_2 = (5,9,4,1)$ --- $v_0$ itself is
untouched by either update, and all three versions answer range-min
queries independently, exactly as the reference implementations
demonstrate.

Separately, for values $(3,1,4,1,5,9,2,6,5,3)$, the range $[3,8]$
(1-indexed, values $4,1,5,9,2,6$) has $3$ elements $\le 4$ and its $3$rd
smallest value is $4$ --- both computed purely from the two chained
versions $S_8$ and $S_2$, cross-checked against brute force over every
possible $(l, r, x)$ and $(l, r, k)$ in the reference code.
\end{examplebox}

\section{Further Reading}
\begin{itemize}[leftmargin=1.5em]
  \item AtCoder Algorithm Lectures index (source of this curriculum):
        \url{https://info.atcoder.jp/entry/algorithm_lectures/index}
  \item J. R. Driscoll, N. Sarnak, D. D. Sleator, R. E. Tarjan, ``Making
        Data Structures Persistent'', STOC 1986, pp.\ 109--121.
  \item USACO Guide, \emph{Persistent Data Structures}:
        \url{https://usaco.guide/adv/persistent?lang=cpp}
\end{itemize}

\end{document}
